\section{Notation, architecture, and training}
\label{app:training}
This appendix defines the architecture, training objectives, and notation. The sensory sequence
clock is indexed by $\tau$; the robot's macro-decision clock is indexed by $n$.
One macro-decision starts at environment step $t_n$ and has nominal horizon
$H=8$. In the closed-loop runner, the transformed command is held for that
interval, unless the episode terminates or reaches its remaining-step limit.

\subsection{Notation at a glance}
\begin{longtable}{@{}p{.21\linewidth}p{.27\linewidth}p{.46\linewidth}@{}}
\caption{Principal symbols. Scalars, dimensions, learned quantities, and fixed
offline quantities are identified separately.}\label{tab:notation}\\
\toprule Symbol & Shape or value & Meaning \\\midrule\endfirsthead
\toprule Symbol & Shape or value & Meaning \\\midrule\endhead
$m,\tau,T,D_m$ & Stream; position; $16$; stream-specific & Sensory stream, sequence index, length, and channel count.\\
$x^m_\tau,\hat x^m_{\tau+1}$ & $\R^{D_m}$ & Observed current and predicted next feature slice.\\
$p_m$ & Positive scalar & Persistence MSE used to normalize the sensory objective.\\
$I_\tau,v_\tau,u_\tau$ & Per-neuron values & Input current, pre-reset membrane, and post-reset membrane.\\
$s_\tau,\beta,V_{\rm th}$ & Binary; $0.85$; $1$ & Spike, decay, and firing threshold.\\
$n,t_n,H$ & Decision; environment step; $8$ & Separate macro-decision and simulator clocks.\\
$q_n,q_n^+,D_q$ & $\R^{39}$; $39$ & Current/outcome physical observations and their dimension.\\
$a_n^*,\tilde a_n,a_n$ & $[-1,1]^4$ for valid commands & Desired, inverse-proposed, and executed commands.\\
$c$ & Task identifier & One of the four manipulation task embeddings.\\
$d_n,\hat d_n$ & $\R^{39}$ & Standardized observed and predicted observation change.\\
$\mu_\Delta,\sigma_\Delta$ & $\R^{39}$, fixed & Offline mean and positive scale of observation changes.\\
$\theta,\theta_0,g_\theta$ & Trained parameters / predictor & Offline parameters; frozen deployment copy; endpoint predictor.\\
$h_n,\tilde h_n,d$ & $\R^{96}$; $96$ & Original/corrected action feature and hidden width.\\
$\bar h,B,r,z_n$ & $\R^{96}$; $\R^{30\times96}$; $30$; $\R^{30}$ & Fixed mean, fixed basis, basis rank, and compressed feature.\\
$\rho,C_\rho$ & Route; $\R^{96\times30}$ & Inferred mechanism and its online residual matrix.\\
$W_{\rm out},b_{\rm out}$ & $\R^{39\times96}$; $\R^{39}$ & Frozen linear endpoint readout.\\
$j,J,T_j,S_j$ & Index; $32$; map; scalar & Candidate, bank size, command transform, and cumulative loss.\\
$k,g,\hat k,\hat g$ & Scalars & Shear/gain and their selected candidate estimates.\\
$\eta,R_C,\delta$ & $0.02$, $0.20$, $0.05$ & Matrix step length, matrix norm bound, command correction cap.\\
$\lambda_b,\lambda_a$ & $0.25$, $4$ & Fixed binding/action loss weights in joint training.\\
$E_M,G_M$ & MSE; percentage & Error of method $M$ and gain relative to frozen prediction.\\
$\ell_m,\mathcal L_{\rm next}$ & Scalar losses & Per-stream and grouped sensory prediction objectives.\\
$\mathcal L_{\rm sem},\mathcal L_{\rm bind},\mathcal L_{\rm act}$ & Scalar losses & Semantic classification, image/text association, and action prediction objectives.\\
$\ell_{n,j}$ & Scalar & Observed transition loss for candidate $j$.\\
$L_n,A_n,C'$ & Scalar; $\R^{96\times30}$; same & Matrix-write loss, its gradient, and tentative matrix before projection.\\
$\epsilon$ & $10^{-12}$ & Safeguard for norm division in the matrix update.\\
$e_n,v_n$ & $\R^{39}$ & Uncorrected prediction residual and output-space correction.\\
$N,y_i,p_{i,y_i}$ & Count; class; probability & Batch size, class label and its predicted probability.\\
$u_i,v_j$ & $\R^{64}$, unit norm & Image and text binding vectors (distinct from membrane variables).\\
$P_i,b_{ij},\tau_b$ & Index set; scalar; $0.12$ & Matching-class texts, image/text similarity and binding temperature.\\
\bottomrule
\end{longtable}

\subsection{Data interfaces and the shared core}
The audio streams use 224 channels for SHD and 285 for SSC. SHD input is
zero-padded to the shared audio stem's width, while the corresponding output
heads retain their original dimensions. Retinal image and rendered-word
streams each use 40 channels and share an input stem. The video stream uses
896 pooled feature channels. All sensory sequences have length 16. Retinal
data come from the digit-image source in \texttt{sklearn.datasets.load\_digits};
rendered words are represented by spatial features under a common scanpath.
Video inputs are causal feature changes from a content-disjoint
Something-Something-V2 cache. The study uses cached inputs with their recorded
splits rather than collecting physical event-camera measurements.

The model has width 96, two shared blocks, four attention heads, and a
feed-forward width of 192. Input stems apply layer normalization and linear
projection, followed by position and modality embeddings. Input spikes are
added to the projected numerical features. Each shared block has spike-based
attention and a spiking feed-forward path, both with residual connections.
The output is layer-normalized. Sensory next-slice prediction adds the learned
output-head change to the current slice, so copying the current feature is
the natural zero-change reference.

The attention implementation restricts each query to a local causal window
of at most four positions, rather than selecting the four largest scores
from all past positions. Queries, keys, and values are formed with LIF units;
membrane-dependent gates modulate attention scores, which are normalized by
softmax within the available window. This limits the number of attention pairs. We do not measure neuromorphic
energy in these experiments. LIF state is initialized
for each input sequence; deployment memory resides in the explicit external
state rather than an unbounded cross-episode membrane history.

The action branch embeds the 39-coordinate observation and four-coordinate
command into two tokens. Task and token-type embeddings are added to them,
and the same core processes the pair. Its final token provides the action
feature $h_n$ for the linear readout. The complete checkpoint has 1,451,388
parameters. The joint-training configuration exposes 605,079 trainable
parameters; the main deployment interface freezes the selected checkpoint.

\subsection{Full offline objectives}
\label{app:losses}
Let $\ell_{\rm SHD}$ and $\ell_{\rm SSC}$ be Equation~\eqref{eq:next} on the
two audio sources, and let $\ell_{\rm image},\ell_{\rm text},\ell_{\rm video}$
be the other sensory losses. The grouping used in training is
\begin{equation}
 \mathcal L_{\mathrm{next}}=\frac14\left[
 \frac{\ell_{\rm SHD}+\ell_{\rm SSC}}2+
 \ell_{\rm image}+\ell_{\rm text}+\ell_{\rm video}\right].
\end{equation}
Audio receives one group weight despite having two sources. The persistence
denominator is floored at $10^{-8}$ in code to prevent division by zero.
This equation defines the training loss, whose denominators come from the
training split. Evaluation aggregation is specified separately in
Appendix~\ref{app:joint}.

For a batch of $N$ labeled examples, let $y_i$ be the class of example $i$
and let $p_{i,y_i}$ be its predicted softmax probability. Semantic cross-entropy
is $-N^{-1}\sum_{i=1}^N\log p_{i,y_i}$, averaged across the five sensory
streams. The coupled heads flatten the length-16 hidden sequence, normalize
it, apply dropout during training, and predict the relevant classes. The
class counts are 20 for SHD, 35 for SSC, 10 for retinal digits/words, and 42
for the video cache. An inherited stable semantic branch separately operates
on standardized flattened sensory inputs. Both branches are frozen at
deployment; the joint-learning semantic results concern the coupled branch.

For binding, $N$ is the number of paired image/text examples in the batch.
The subscripts $I$ and $T$ below denote image and text. Let $u_i$ and $v_j$ be
unit-normalized 64-dimensional image and
text representations, respectively. These are distinct from the membrane
symbol $u_\tau$ in Equation~\eqref{eq:lif}. Define similarity
$b_{ij}=u_i^\top v_j/\tau_b$, where the temperature is $\tau_b=0.12$.
Let $P_i$ contain all text indices whose digit class matches image $i$.
The image-to-text loss is
\begin{equation}
 \mathcal L_{I\rightarrow T}=-\frac1N\sum_{i=1}^N
 \log\frac{\sum_{j\in P_i}\exp(b_{ij})}{\sum_{j=1}^N\exp(b_{ij})}.
\end{equation}
This assigns probability mass to all matching-class texts, rather than one
arbitrarily chosen pair. The reverse loss $\mathcal L_{T\rightarrow I}$
exchanges image and text roles, and
$\mathcal L_{\mathrm{bind}}=(\mathcal L_{I\rightarrow T}+\mathcal L_{T\rightarrow I})/2$.
The binding representation averages hidden features across sequence positions
before its projection.

The action objective is Equation~\eqref{eq:action}. Its mean and scale are
fit from offline observation changes and reused during deployment. The
weights in Equation~\eqref{eq:joint} are the fixed configuration used by the
recorded study. Sequential action training runs for 800 updates. The additional
joint stage uses 200 updates, batch size 32, and learning rate $10^{-4}$.
Sensory prediction and action prediction both supply gradients to the shared
core in this joint stage.
